\section*{CHIMIE (7Points) : Transformations chimiques spontanées}
Les transformations chimiques diffèrent selon le type de systèmes chimiques et les conditions initiales, et sont soit rapides ou lentes. Certaines d'entre elles conduisent à la synthèse de produits, et peuvent être utilisées dans différents domaines tels que la santé et l'industrie, et ce selon des protocoles déterminés.\\
Cet exercice vise, l'étude de la méthode de contrôle de l'évolution d'un système chimique à partir de réaction de synthèse de l'aspirine (acide acétylsalicylique), et l'étude du comportement des molécules de cet acide dans l'eau afin de déterminer sa constante d'acidité, ainsi que l'étude d'une transformation spontanée dans une pile.\\

\section*{Première partie : synthèse de l'aspirine au laboratoire, et étude de sa réaction avec l'eau}
\begin{enumerate}
  \item L'acide acétylsalicylique ou aspirine peut être synthétisé au laboratoire à partir de la réaction entre l'acide salicylique et l'anhydride éthanoïque en utilisant le chauffage à reflux selon l'équation de la réaction suivante modélisant cette transformation :\\
\includegraphics[width=\textwidth]{2025_10_31_838bdd6d133f8b78fc82g-1}\\
1.1. Donner le nom du groupement fonctionnel délimité par un trait pointillé fermé dans la forme topologique de chacune des molécules d'acide salicylique et d'acide acétylsalicylique.\\
1.2. Citer les deux caractéristiques de cette transformation.\\
1.3. Choisir, parmi les montages expérimentaux (1), (2) et (3) le montage utilisé pour réaliser cette synthèse.\\
\includegraphics[width=\textwidth]{2025_10_31_838bdd6d133f8b78fc82g-1(1)}\\
1.4. Quel est l'intérêt du chauffage à reflux ?\\
1.5. On introduit dans une fiole jaugée, $n_{1}=0,10 \mathrm{~mol}$ d'acide salicylique et $n_{2}=0,26 \mathrm{~mol}$ d'anhydride éthanoïque et quelques gouttes d'acide sulfurique concentré. Après chauffage à reflux, et les opérations de traitement et de purification, on obtient des cristaux d'aspirine de masse $m_{\text {exp }}=15,3 \mathrm{~g}$.\\
Calculer le rendement de cette synthèse sachant que le réactif limitant est l'acide salicylique.\\
On donne : Masse molaire de l'acide acétylsalicylique : $M=180 \mathrm{~g} \cdot \mathrm{~mol}^{-1}$
  \item On prépare une solution aqueuse ( $S$ ) d'acide acétylsalicylique de concentration molaire\\
$C=5,55.10^{-3} \mathrm{~mol} . \mathrm{L}^{-1}$ et de volume $V=500 \mathrm{~mL}$. Après mesure de la conductivité de la solution $(S)$, on détermine la valeur de $x_{\mathrm{f}}$ avancement de la réaction à l'état final du système chimique, et on trouve: $x_{\mathrm{f}}=5,70 \cdot 10^{-4} \mathrm{~mol}$.\\
Pour simplifier, on désigne la molécule de l'acide acétylsalicylique par $A H$, et sa base conjuguée par $A^{-}$.\\
2.1. Ecrire l'équation de la réaction de l'acide acétylsalicylique avec l'eau.\\
2.2. Montrer que la réaction de l'acide acétylsalicylique avec l'eau est non totale.\\
2.3. Déterminer la valeur de la constante d'acidité $K_{\mathrm{A}}$ du couple $A H_{(\mathrm{aq})} / A_{(\mathrm{aq})}^{-}$.
\end{enumerate}



\section*{Deuxième partie : Transformation spontanée dans une pile}
On réalise une pile en utilisant le matériel et les produits suivants~:
\begin{itemize}
    \item un bêcher contenant le volume $V_{1}=\qty{20}{\mL}$ d'une solution
          aqueuse de nitrate d'argent $\ce{Ag^+(aq) + NO^-_3_{(aq)}}$ de
          concentration molaire $C_{1}=\qty{1.0e-1}{\mol\per\L}$
    \item un bêcher contenant le volume $V_{2}=\qty{20}{\mL}$ d'une solution
          aqueuse de nitrate de cuivre $\ce{Cu^2+_{(aq)} + 2NO^-_3_{(aq)}}$ de
          concentration molaire $C_{2}=\qty{5.0e-2}{\mol\per\L}$
    \item un fil de cuivre~;
    \item un fil d'argent~;
    \item un pont salin contenant une solution aqueuse saturée de nitrate de
          potassium $\ce{K^+(aq) + NO^-_3(aq)}$.
\end{itemize}

\begin{donnees}
    \begin{itemize}
        \item $1\mathcal{F}=\qty{96500}{\C\mol}$~;
        \item Constante d'équilibre associée à l'équation
              $\ce{2Ag^{+}_{(aq)} + Cu_{(s)} <=>[(1)][(2)] 2Ag_{(s)} +
              Cu^{2+}_{(aq)}}$ est~: $\mathrm{K}=\num{2.2e15}$
    \end{itemize}
\end{donnees}

On relie les électrodes de la pile à un conducteur ohmique en série avec un
ampèremètre, et on observe le passage d'un courant électrique dans le circuit
extérieur de la pile.
\begin{enumerate}
    \item Calculer la valeur du quotient de la réaction $Q_{r,i}$ dans l'état
          initial du système chimique. En déduire le sens spontané de
          l'évolution de ce système.
    \item On fait fonctionner la pile pendant une longue durée jusqu'à ce qu'il
          s'épuise. Déterminer la valeur de la quantité d'électricité qui
          traverse le conducteur ohmique depuis le début de fonctionnement de la
          pile jusqu'à son épuisement, sachant que le réactif limitant est l'ion
          $\ce{Ag^+}$.
\end{enumerate}




